% ============================================================
%  2026 国赛论文骨架（由 cumcmthesis 官方模板派生）
% ------------------------------------------------------------
%  使用说明：
%  1. 把这个文件上传到 Overleaf 项目里，在 Menu 里设为 Main document
%  2. 编译器务必选 XeLaTeX
%  3. 电子版提交：保持 withoutpreface 选项（不含承诺书和编号页）
%     需要打印纸质版时：改成 \documentclass[bwprint]{cumcmthesis}
%  4. 全文任何位置（含页眉、图表、附录）都不得出现校名、姓名、学号
%  5. 正文（不含附录）不超过 30 页
% ============================================================

\documentclass[withoutpreface,bwprint]{cumcmthesis}

\usepackage[most]{tcolorbox}
\usepackage{tikz}
\usepackage{cumcm2026}
\usepackage{url}
\usepackage{subcaption}

\title{论文题目（定题后改这里：对象 + 方法 + 问题，20 字以内）}

% 注意：题号、报名号、学校、队员姓名等命令已废弃，不要填！填了会违规。

\begin{document}

\maketitle

% ---------- 摘要：最后一天上午专门磨 2 小时，600~900 字，纯文字无公式图表 ----------
\begin{abstract}
% 总起段：本文针对……背景下的 X、Y、Z 问题，基于……理论，分别建立了……模型，并通过……算法求解。
%
% 每问一段（赛题分几问就几段）：针对问题一，考虑到……，建立了……模型；采用……方法求解，
% 得到……（必须给具体数字）。结果表明……（定量结论）。
%
% 收尾段：最后对模型进行了灵敏度分析与误差检验，结果表明模型……；模型可推广至……场景。
%
% 关键词 3~5 个，含具体方法名（如 熵权法、TOPSIS、遗传算法）
\end{abstract}

\section{问题重述}
% 用自己的语言复述赛题背景与要解决的问题，不要照抄原题。0.5~1 页。
% 建议分两块：「问题背景」+「问题提出/要求」。

\section{问题分析}
% 展示建模思路与选模型的依据——这部分最能体现思考过程。
% \subsection{问题一的分析}
% \subsection{问题二的分析}
% \subsection{问题三的分析}

\section{模型假设与符号说明}
\subsection{模型假设}
% 3~7 条，每条都要有依据，不要罗列废话。
% 示例：假设车辆匀速行驶。——依据：附件数据显示 95% 行程速度方差小于 5 km/h。

\subsection{符号说明}
% 三线表，符号 + 含义 + 单位，全文统一（见文末语法速查）。

\section{数据预处理与模型准备}
% 可选章节。数据来源说明、缺失值/异常值处理、标准化方法。
% 注意：异常值可能是关键信息，别无脑剔除。

\section{问题一的模型建立与求解}
% 固定结构：模型建立 → 求解方法（配算法流程图）→ 结果分析（配图表 + 定量结论）

\section{问题二的模型建立与求解}

\section{问题三的模型建立与求解}

\section{问题四的模型建立与求解}
% 按赛题实际问数增删，用不上直接删掉本节省版面。

\section{模型检验与灵敏度分析}
% 国奖论文标配：误差分析 + 灵敏度分析（参数扰动后结果是否稳定）。

\section{模型评价与推广}
\subsection{模型的优点}
\subsection{模型的不足与改进方向}
% 诚实陈述，缺点别藏着；给出适用范围和可推广的场景。

% ============================================================
%  AI 使用声明（2026 新规：放在参考文献之前，二选一）
%  用了 AI：改下面的用途文字；没用 AI：把它换成 \CumcmAINotUsed
% ============================================================
\CumcmAIUsed{语言润色、\LaTeX{} 排版辅助与代码调试}
% \CumcmAINotUsed

% ============================================================
%  参考文献：GB/T 7714 格式，正文引用与此处一一对应
% ============================================================
\begin{thebibliography}{9}
  \bibitem{ref1} 作者. 题名[文献类型标识]. 刊名, 年, 卷(期): 起止页码.
  \bibitem{ref2} 作者. 题名[文献类型标识]. 刊名, 年, 卷(期): 起止页码.
  % 例：刘海洋. \LaTeX{} 入门[M]. 北京: 电子工业出版社, 2013.
\end{thebibliography}

% ============================================================
%  附录：支撑材料文件列表 + 全部可运行源代码（页数不限）
%  注意：附录参与查重，不要放 AI 生成的文字内容
% ============================================================
\begin{appendices}
  \section{支撑材料文件列表}
  % 列出支撑材料 ZIP 里的每个文件及其作用（官方要求必须列入）。

  \section{程序代码}
  % 用 lstlisting 环境贴全部可运行代码，缺失或与结果不符可能被取消评奖资格。
  % \begin{lstlisting}[language=Python]
  % import numpy as np
  % \end{lstlisting}
\end{appendices}

% ============================================================
%  语法速查（需要时解除注释；不需要就一直注释着，不影响编译）
% ============================================================
% 插图：
% \begin{figure}[!htbp]
%   \centering
%   \includegraphics[width=0.75\textwidth]{figures/文件名.png}
%   \caption{图名：要自明，带单位}
%   \label{fig:xxx}       % 文中用 如图~\ref{fig:xxx} 所示 引用
% \end{figure}
%
% 三线表：
% \begin{table}[!htbp]
%   \centering
%   \caption{表名}
%   \begin{tabular}{ccc}
%     \toprule[1.5pt]
%     符号 & 含义 & 单位 \\
%     \midrule[1pt]
%     v    & 速度 & m/s  \\
%     t    & 时间 & s    \\
%     \bottomrule[1.5pt]
%   \end{tabular}
% \end{table}
%
% 带编号公式（行间，可引用）：
% \begin{equation}
%   E = mc^{2}
%   \label{eq:energy}
% \end{equation}
% 行内公式：$E=mc^2$；无编号行间公式：\[ E=mc^2 \]
% 多行对齐用 align 环境，& 为对齐位置。

\end{document}
